\documentclass[12pt]{article}
\usepackage[utf8]{inputenc}
\usepackage[T1]{fontenc}
\usepackage{hyperref}
\title{Enrekang Language Resources for K-12 Curriculum Development}
\author{Compiled by Claude Sonnet 5.5 (claude-sonnet-5-5) for Andr\'as Kornai}
\date{}
\begin{document}
\maketitle

\section{Introduction}
Enrekang (ISO 639-3 ptt; Glottocode enre1239) is an Austronesian language of the Northern South Sulawesi branch, spoken in Enrekang Regency, Indonesia. It is closely related to Duri, Maiwa and Malimpung (the Massenrempulu cluster). Searches found only general descriptions and a few word lists. No Enrekang primer, textbook, dictionary or grammar turned up, so the material below is thin.

\section{Primary Educational Resources}
No primers, textbooks or school programmes for Enrekang were found in the search results.

\section{Linguistic and Reference Materials}
\begin{itemize}
\item Wikipedia, ``Enrekang language'' (a stub giving classification and a speaker estimate): \url{https://en.wikipedia.org/wiki/Enrekang_language}
\item ASJP word list for Enrekang: \url{https://asjp.clld.org/languages/ENREKANG}
\item Omniglot, Duri (a closely related Massenrempulu language; the search result describes the cluster and dialects): \url{https://omniglot.com/writing/duri.htm}
\item Duri dictionary project on Webonary (for the related language Duri, not Enrekang itself): \url{https://www.webonary.work/duri/?p=9444}
\end{itemize}

\section{Religious and Community Resources}
None found in the search results.

\section{Digital and Audio Resources}
None specific to Enrekang were found. The Duri Webonary site above may serve as a model or a source of cognate material.

\section{Curriculum Development Context}
The search results gave no evidence of Enrekang being taught in schools. A curriculum would first need basic documentation (word lists, texts, orthography). Related Duri materials may help, but they should not be assumed to apply to Enrekang without checking with speakers.
\end{document}